% ============================================================================
% INFORME DE CIERRE — PID PP9884
% Análisis del Proceso de Liofilizado Aplicado a Alimentos
% de Producción Local/Regional en la Provincia de Tierra del Fuego
% UTN FRTDF — Grupo QA3DS
% ============================================================================
\documentclass[12pt, a4paper]{article}

% --- Idioma y codificación ---
\usepackage[utf8]{inputenc}
\usepackage[T1]{fontenc}
\usepackage[spanish,es-tabla,es-noshorthands]{babel}

% --- Diseño de página ---
\usepackage[margin=2.5cm, headheight=15pt]{geometry}
\usepackage{fancyhdr}
\usepackage{lastpage}

% --- Tipografía y colores ---
\usepackage{xcolor}
\usepackage{lmodern}

% --- Tablas y gráficos ---
\usepackage{booktabs}
\usepackage{longtable}
\usepackage{array}
\usepackage{graphicx}
\graphicspath{{../figuras/}}
\usepackage{float}

% --- Otros ---
\usepackage{enumitem}
\usepackage{hyperref}
\hypersetup{
    colorlinks  = true,
    linkcolor   = blue!60!black,
    urlcolor    = blue!50!black,
    citecolor   = green!50!black,
}
\usepackage{multirow}

% --- Comando para porcentaje (evita conflicto babel) ---
\newcommand{\pct}{\char`\%}

% --- Encabezado y pie ---
\pagestyle{fancy}
\fancyhf{}
\fancyhead[L]{\small PID PP9884 --- Informe de Cierre}
\fancyhead[R]{\small UTN FRTDF --- QA3DS}
\fancyfoot[C]{\small Página \thepage\ de \pageref{LastPage}}

% ============================================================================
\begin{document}

% --- PORTADA ---
\begin{titlepage}
\centering
\vspace*{2cm}

{\Large\bfseries Universidad Tecnológica Nacional}\\[0.3cm]
{\large Facultad Regional Tierra del Fuego}\\[0.3cm]
{\large Secretaría de Ciencia, Tecnología y Posgrado}\\[1.5cm]

\rule{\linewidth}{1pt}\\[0.5cm]
{\LARGE\bfseries INFORME DE CIERRE}\\[0.3cm]
{\Large Proyecto de Investigación y Desarrollo}\\[0.3cm]
\rule{\linewidth}{1pt}\\[1cm]

{\Large\bfseries Código: PP9884}\\[0.5cm]
{\large\itshape Análisis del Proceso de Liofilizado Aplicado a Alimentos\\
de Producción Local/Regional en la Provincia de Tierra del Fuego}\\[1.5cm]

{\large Grupo de Investigación:}\\[0.2cm]
{\Large\bfseries QA3DS}\\[0.2cm]
{\normalsize Química Aplicada al Ambiente, los Alimentos y el Desarrollo Sostenible}\\[1.5cm]

\begin{tabular}{ll}
\textbf{Director:} & Ing. Pesq. Ariel Luján Giamportone \\
\textbf{Período:} & 1 de abril de 2023 --- 31 de marzo de 2025 \\
\textbf{Prórroga:} & DISPOSICIÓN-1-2025-PRORROGA-PP9884 \\
\textbf{Sede:} & Ushuaia / Río Grande, Tierra del Fuego \\
\end{tabular}

\vfill
{\normalsize Documento preliminar --- Marzo 2026}

\end{titlepage}

% --- TABLA DE CONTENIDOS ---
\tableofcontents
\newpage

% ============================================================================
\section{Datos generales del proyecto}
% ============================================================================

\begin{longtable}{p{5cm}p{10cm}}
\toprule
\textbf{Campo} & \textbf{Detalle} \\
\midrule
Código & PP9884 (PIDPP UTN) \\
Título & Análisis del Proceso de Liofilizado Aplicado a Alimentos de Producción Local/Regional en la Provincia de Tierra del Fuego \\
Institución & Universidad Tecnológica Nacional --- FRTDF \\
Grupo & QA3DS (Química Aplicada al Ambiente, los Alimentos y el Desarrollo Sostenible) \\
Director & Ing. Pesq. Ariel Luján Giamportone \\
Período original & 1/04/2023 al 31/03/2025 \\
Prórroga & Aprobada por DISPOSICIÓN-1-2025 (15/03/2025) \\
Sede experimental & Laboratorio UTN FRTDF (Río Grande) \\
Financiamiento & Sistema PIDPP --- UTN \\
\bottomrule
\end{longtable}

% ============================================================================
\section{Equipo de investigación}
% ============================================================================

\subsection{Investigadores docentes}

\begin{longtable}{p{5.5cm}p{4cm}p{4cm}}
\toprule
\textbf{Nombre} & \textbf{Título} & \textbf{Rol} \\
\midrule
Ariel Luján Giamportone & Ing. Pesquero & Director \\
Pamela Rocío Flores & Ing. Química & Investigadora \\
María Victoria Cornejo & Ing. Química & Investigadora \\
Javier Alfarano & Ing. Industrial & Investigador de Apoyo \\
Milagro Mottola & Dra. en Química & Investigadora \\
Tomás Chalde & Dr. & Investigador de Apoyo \\
\bottomrule
\end{longtable}

\subsection{Alumnos investigadores}

\begin{longtable}{p{4.5cm}p{3cm}p{2cm}p{2cm}p{2cm}}
\toprule
\textbf{Nombre} & \textbf{Carrera} & \textbf{Ingreso} & \textbf{Egreso} & \textbf{Beca} \\
\midrule
\multicolumn{5}{l}{\textit{Activos al cierre}} \\
\midrule
Emanuel Silva Naveas & Electromecánica & Jul 2024 & --- & BINID 2025 \\
Gabriel F. Costilla & Pesquera & Abr 2023 & --- & GyJ 2025 \\
Jorge Rojas Flores & --- & Abr 2023 & --- & GyJ 2025 \\
Angélica Cárcamo & Ing. Química (grad.) & Mar 2025 & --- & BINID 2025 \\
Ricardo A. Heredia & Electromecánica & Jul 2024 & --- & BINID 2025 \\
Facundo N. Gutiérrez & Ing. Química & Abr 2024 & --- & --- \\
Matías I. Álvarez & Ing. Química & Abr 2024 & --- & BINID 2025 \\
\midrule
\multicolumn{5}{l}{\textit{Egresados}} \\
\midrule
Michael A. Trujillo Burgos & Química & Abr 2023 & Mar 2025 & --- \\
Federico D. Cejas & --- & Abr 2023 & Mar 2024 & --- \\
Roxana B. Barrionuevo & --- & Abr 2023 & Mar 2024 & --- \\
Franco Guereta & --- & Abr 2023 & Mar 2024 & --- \\
Erika G. Mamaní Mejía & --- & Abr 2023 & Mar 2024 & --- \\
Vanina Chávez & Química & Abr 2024 & --- & --- \\
Facundo Fernández B. & Química & Abr 2024 & --- & --- \\
Tamara P. Godoy Ramos & --- & Abr 2023 & --- & --- \\
Daniel M. Esteban & Pesquera & Abr 2023 & --- & --- \\
\bottomrule
\end{longtable}

\textbf{Resumen de Recursos Humanos formados:} 6 docentes investigadores + 16 alumnos participaron a lo largo del proyecto. Al cierre se mantienen activos 7 alumnos con dedicación de 5\,h/semana. Se gestionaron 6 becas (4 BINID 2025, 2 becas ``Gabriel y Jorge'' 2025).

% ============================================================================
\section{Objetivos del proyecto: cumplimiento}
% ============================================================================

\subsection{Objetivo general}

\textit{``Agregar valor a productos regionales abriendo nuevos mercados y generando oportunidades para productores locales mediante una técnica innovadora (liofilización), mejorando la cadena de suministro con características de soberanía alimentaria en Tierra del Fuego.''}

\textbf{Evaluación:} Se logró avance significativo en la generación de conocimiento técnico y científico sobre liofilización aplicada a ruibarbo patagónico. Se establecieron parámetros óptimos de proceso, se completó un estudio de escalado económico y se inició la vinculación con el sector productivo (Estancia Viamonte). Las matrices de salicornia y erizo de mar fueron descartadas por limitaciones logísticas y de acceso a materia prima.

\subsection{Objetivos específicos}

\begin{longtable}{p{0.5cm}p{8cm}p{2cm}p{3cm}}
\toprule
\textbf{\#} & \textbf{Objetivo} & \textbf{Estado} & \textbf{Observación} \\
\midrule
OE1 & Generar know-how sobre productos de flora y fauna típica de TDF con potencial comercial & Parcial & Solo ruibarbo (salicornia y erizo descartados) \\
OE2 & Explorar y mejorar la obtención de alimentos con valor agregado mediante liofilización & Cumplido & Protocolo optimizado para ruibarbo \\
OE3 & Optimizar parámetros de liofilización para mejorar calidad y estabilidad & Cumplido & Modelo de Page ajustado, plateau 36--48\,h \\
OE4 & Evaluar métodos de análisis de datos para caracterizar productos liofilizados & Cumplido & EDA, ANOVA II, Page, Kruskal-Wallis implementados \\
OE5 & Identificar factores clave que influyen en la calidad y estabilidad & Parcial & Pretratamiento y tiempo identificados; faltan análisis bioquímicos y sensoriales \\
OE6 & Desarrollar equipamiento técnico: diseño de liofilizador local & Pendiente & Se formuló estudio de escalado; diseño propio no se inició \\
OE7 & Transferir conocimiento y tecnología al sector socioproductivo & En proceso & Convenio Estancia Viamonte activo, artículos en preparación \\
OE8 & Consolidar el grupo de I+D+I QA3DS en UTN FRTDF & Cumplido & Grupo operativo, presencia digital, 6 becados \\
\bottomrule
\end{longtable}

\textbf{Síntesis:} De los 8 objetivos específicos, 4 se cumplieron satisfactoriamente, 2 se cumplieron parcialmente (por la reducción de matrices de 3 a 1), 1 está en proceso y 1 quedó pendiente (no era alcanzable en el período del PID sin financiamiento adicional).

% ============================================================================
\section{Actividades realizadas}
% ============================================================================

\subsection{Trabajo experimental}

Se realizaron cuatro baterías experimentales de liofilización de ruibarbo (\textit{Rheum rhabarbarum} L.) entre diciembre de 2023 y enero de 2025:

\begin{longtable}{p{1.5cm}p{2cm}p{2.5cm}p{2cm}p{5cm}}
\toprule
\textbf{Exp.} & \textbf{Fecha} & \textbf{Pretratamiento} & \textbf{Réplicas} & \textbf{Resultado principal} \\
\midrule
Exp1 & Dic 2023 & FRESCO & 3 & Primera curva cinética; pérdida $\sim$94\,\pct{} BH \\
Exp2 & May 2024 & FRESCO, CONG., ULTRACONG. & 9 & Primer diseño factorial 3$\times$3 \\
Exp3 & Jul 2024 & FRESCO, CONG., ULTRACONG. & 9 & Serie más completa; plateau a 36\,h \\
Exp4 & Ene 2025 & FRESCO, CONG., ULTRACONG. & 9 & Confirmación de tendencias \\
\bottomrule
\end{longtable}

\textbf{Dataset total:} 261 observaciones (30 combinaciones pretratamiento\,$\times$\,tiempo, con 9 réplicas cada una).

\textbf{Equipamiento utilizado:}
\begin{itemize}[nosep]
    \item Liofilizador RIFICOR LT-8 (laboratorio UTN FRTDF, Río Grande)
    \item Balanza analítica OHAUS AR2140 (210\,g, 0{,}1\,mg)
    \item Balanza portátil OHAUS CR (2200\,g, 1\,g) --- adquirida jun 2024
    \item Termómetro digital Alla France $-50$ a $+70$\,°C --- adquirido jun 2024
    \item Selladora al vacío + bolsas/rollos gofrados Turbosaver
\end{itemize}

\textbf{Muestras:} Peciolos de ruibarbo cortados transversalmente ($\sim$1\,cm) provenientes de huerta doméstica (Río Grande) y Estancia Viamonte (30\,km de Río Grande). Dos ecotipos: invernadero (más turgente) y exterior (más blando).

\subsection{Resultados principales}

Los resultados experimentales completos se presentan en el preprint científico adjunto (\texttt{main.tex}). Se resumen a continuación:

\begin{enumerate}[nosep]
    \item \textbf{Cinética de secado:} El modelo de Page describe adecuadamente la pérdida de peso ($R^2 = 0{,}89$--$0{,}92$). El pretratamiento ULTRACONGELADO mostró la cinética más rápida ($k = 0{,}1067$\,h$^{-n}$, $n = 0{,}8323$).

    \item \textbf{Tiempo óptimo:} Se alcanza pérdida de peso del 94\,\pct{} a las 36\,h con barquillo de aluminio. El plateau de la derivada numérica confirma estabilización:
    \begin{itemize}[nosep]
        \item FRESCO: 72\,h
        \item CONGELADO: 48\,h
        \item ULTRACONGELADO: 48\,h
    \end{itemize}

    \item \textbf{Efecto del pretratamiento:} Significativo a las 24\,h (Kruskal-Wallis $p = 0{,}0004$), pero convergente a 36--48\,h ($p > 0{,}05$). Implicancia industrial: el congelamiento previo acelera las primeras fases pero no altera el resultado final.

    \item \textbf{Reproducibilidad:} La ANOVA Tipo~II ($R^2 = 0{,}923$) mostró que el factor mes (bloque) no es significativo ($p = 0{,}1005$), confirmando que los resultados son reproducibles entre baterías experimentales.

    \item \textbf{Condiciones operacionales:} Temperatura del condensador $-36$ a $-41$\,°C; presión 0{,}94 a 2{,}432\,mmHg. El equipo no siempre alcanzó las especificaciones mínimas del manual.
\end{enumerate}

\subsection{Matrices descartadas}

\begin{itemize}[nosep]
    \item \textbf{Salicornia (\textit{Sarcocornia magellanica}):} Descartada por dificultades de acceso a materia prima en cantidades suficientes y por la falta de protocolo previo de procesamiento. Se conserva la línea para un futuro proyecto.
    \item \textbf{Erizo de mar (\textit{Loxechinus albus}, \textit{Pseudechinus magellanicus}):} Descartado por complejidad logística de obtención de gónadas frescas en condiciones controladas, restricciones de pesca y necesidad de protocolos de bioseguridad adicionales. Se recopiló documentación bibliográfica (ERIZOS v2--v4.pptx).
\end{itemize}

% ============================================================================
\section{Producción científica y académica}
% ============================================================================

\subsection{Publicaciones y documentos técnicos}

\begin{longtable}{p{0.5cm}p{7.5cm}p{2.5cm}p{2.5cm}}
\toprule
\textbf{\#} & \textbf{Producto} & \textbf{Tipo} & \textbf{Estado} \\
\midrule
1 & ``Química Aplicada al Ambiente, los Alimentos y el Desarrollo Sostenible: Investigando la Liofilización en Tierra del Fuego'' (Flores, Mottola, Cornejo, Alfarano, Giamportone) & Artículo divulg. (La Lupa) & Redactado \\
2 & ``Liofilización: Conservando el Futuro Alimentario de Tierra del Fuego'' & Artículo divulg. (La Lupa) & Redactado \\
3 & Artículo breve La Lupa 2025 (Alfarano) & Artículo divulg. & En preparación \\
4 & Preprint cinética de liofilización de ruibarbo (15\,pp, 7 tablas, 3 figuras) & Preprint científico & Completado \\
5 & Estudio de escalado y factibilidad (5 documentos) & Informe técnico & Completado \\
6 & Business Model Canvas (Mottola) & Análisis de negocio & Completado \\
7 & Guía de uso rápido del liofilizador v1 & Manual técnico & Completado \\
8 & Notebook de análisis estadístico compilado (25 celdas, 6 figuras) & Análisis computac. & Completado \\
9 & Dashboard integral del proyecto (Dash + SQLite, 4 pestañas) & Herramienta web & Completado \\
\bottomrule
\end{longtable}

\subsection{Presentaciones y eventos}

\begin{longtable}{p{3cm}p{7cm}p{3.5cm}}
\toprule
\textbf{Fecha} & \textbf{Evento} & \textbf{Tipo} \\
\midrule
Oct 2023 & Jornadas de Ciencia y Tecnología 2023, UTN FRTDF & Presentación \\
Oct 2024 & Capacitación uso del liofilizador (4\,h presencial) & Capacitación interna \\
2025 & Semana del Ambiente, Municipio de Río Grande & Presentación pública \\
\bottomrule
\end{longtable}

\subsection{Posters}
\begin{itemize}[nosep]
    \item Poster QA3DS PID --- difusión institucional
    \item Poster QA3DS Colegio de Ingenieros
\end{itemize}

\subsection{Infraestructura digital generada}

\begin{itemize}[nosep]
    \item \textbf{GitHub:} Organización \href{https://github.com/QA3DS}{QA3DS}, repositorio \href{https://github.com/QA3DS/PID-Liofilizados}{PID-Liofilizados} (13 issues, 3 milestones, 12 labels), Project Board Kanban
    \item \textbf{LinkedIn:} Página institucional del grupo QA3DS
    \item \textbf{ResearchGate:} Laboratorio QA3DS
\end{itemize}

% ============================================================================
\section{Formación de recursos humanos}
% ============================================================================

\begin{itemize}[nosep]
    \item \textbf{16 alumnos investigadores} participaron a lo largo del proyecto (7 activos al cierre)
    \item \textbf{6 becas gestionadas:} 4 BINID 2025 (Silva Naveas, Cárcamo, Heredia, Álvarez) + 2 becas ``Gabriel y Jorge'' 2025 (Costilla, Rojas Flores)
    \item \textbf{1 capacitación formal:} Uso del liofilizador RIFICOR LT-8 (4\,h, oct 2024), con certificación a participantes. Docentes capacitadores: Cornejo, Flores, Alfarano
    \item \textbf{1 graduada incorporada:} Angélica Cárcamo (Ing. Química) se incorporó como investigadora graduada en mar 2025 con beca BINID
    \item \textbf{Habilidades desarrolladas:} técnicas de liofilización, análisis gravimétrico, procesamiento de datos experimentales, análisis estadístico con Python, gestión de proyectos en GitHub
\end{itemize}

% ============================================================================
\section{Vinculación y transferencia}
% ============================================================================

\begin{itemize}[nosep]
    \item \textbf{Estancia Viamonte:} Convenio de colaboración activo desde 2024 para acceso a materias primas y entorno productivo real (ruibarbo, ajo negro, grosellas)
    \item \textbf{INTI:} Se identificó como referencia de transferencia tecnológica (modelo Trufas del Nuevo Mundo --- trufas liofilizadas exportadas a España y Francia)
    \item \textbf{En gestión:} contactos con INTA, ANMAT, SENASA, Secretaría Provincial de Ambiente
\end{itemize}

% ============================================================================
\section{Ejecución presupuestaria}
% ============================================================================

\subsection{Compras realizadas}

\begin{longtable}{p{5cm}p{2.5cm}p{2cm}p{3.5cm}}
\toprule
\textbf{Ítem} & \textbf{Monto} & \textbf{Fecha} & \textbf{Observación} \\
\midrule
\multicolumn{4}{l}{\textit{Equipamiento (Ítem 4.3)}} \\
\midrule
Balanza portátil OHAUS CR & USD\,230 & Jun 2024 & Prov.: One Lab \\
Termómetro Alla France $-50/+70$°C & USD\,44 & Jun 2024 & Prov.: One Lab \\
\midrule
\multicolumn{4}{l}{\textit{Insumos (Ítem 2)}} \\
\midrule
Rollos gofrados Turbosaver 22\,cm$\times$5\,m ($\times$2) & ARS\,20.260 & Nov 2024 & Envases al vacío \\
Bolsas gofradas Turbosaver 22$\times$30\,cm ($\times$50) & ARS\,13.720 & Nov 2024 & Envases al vacío \\
Bolsas envasado al vacío (compra add.) & ---\textsuperscript{*} & Mar 2025 & Nota PIDPP \\
\midrule
\multicolumn{4}{l}{\textit{Servicios y viáticos (Ítem 3)}} \\
\midrule
Reintegro viáticos dic 2024 & ---\textsuperscript{*} & Dic 2024 & Solicitud V. Cornejo \\
\midrule
\multicolumn{4}{l}{\textit{Solicitudes en gestión}} \\
\midrule
Mesa antivibratoria (para OHAUS AR2140) & ---\textsuperscript{*} & Solicitada & Necesaria para calibración \\
Parrillas/estantes para liofilizador & ---\textsuperscript{*} & Mar 2024 & Solicitada \\
Selladora al vacío & ---\textsuperscript{*} & Solicitada & Prov.: Turbosaver \\
\bottomrule
\end{longtable}

{\footnotesize \textsuperscript{*} Monto no disponible en la documentación digital accesible. Consultar comprobantes físicos.}

\textbf{Nota:} Existen cotizaciones documentadas (N° 56952, N° 49561, Presupuesto 73736, P-00004577) en los archivos del proyecto cuyo detalle de montos ejecutados requiere verificación con los comprobantes originales.

\subsection{Crédito presupuestario disponible}

En fecha 17 de marzo de 2026 se notificó la disponibilidad del crédito correspondiente al saldo al 31/12/2025:

\begin{table}[H]
\centering
\caption{Crédito presupuestario disponible --- PID PP9884 (saldo 31/12/2025)}
\label{tab:credito}
\begin{tabular}{clr}
\toprule
\textbf{Ítem} & \textbf{Concepto} & \textbf{Crédito (ARS)} \\
\midrule
2 & Insumos & \$\,100.577{,}54 \\
3 & Servicios / viáticos / honorarios & \$\,55.498{,}05 \\
4.3 & Equipamiento & \$\,87.485{,}89 \\
\midrule
& \textbf{TOTAL} & \textbf{\$\,243.561{,}48} \\
\bottomrule
\end{tabular}
\end{table}

\subsection{Propuesta estratégica de inversión del saldo}

El saldo disponible de ARS\,243.561{,}48 puede asignarse según el siguiente análisis de prioridades:

\begin{longtable}{p{1.5cm}p{5cm}p{2.5cm}p{4.5cm}}
\toprule
\textbf{Ítem} & \textbf{Propuesta} & \textbf{Monto est.} & \textbf{Justificación} \\
\midrule
\multicolumn{4}{l}{\textit{Prioridad ALTA --- Completar experimentos en curso}} \\
\midrule
2 & Insumos para Exp5+ (ruibarbo fresco, envases, sales para aw, reactivos) & ARS\,$\sim$60.000 & Necesarios para completar el ciclo experimental con análisis de actividad de agua \\
2 & Bolsas y rollos gofrados adicionales & ARS\,$\sim$25.000 & Stock de envasado para experimentos y muestras de conservación \\
3 & Análisis externos: IIByT-CONICET (DVS), CIATI (microbiología, rotulado nutricional) & ARS\,$\sim$55.000 & Dato crítico para OE5 y para el artículo científico D2 \\
\midrule
\multicolumn{4}{l}{\textit{Prioridad MEDIA --- Mejora de infraestructura}} \\
\midrule
4.3 & Mesa antivibratoria para balanza OHAUS AR2140 & ARS\,$\sim$87.000 & Solicitada y pendiente; mejora la precisión analítica \\
\midrule
\multicolumn{4}{l}{\textit{Prioridad BAJA --- Difusión}} \\
\midrule
3 & Viáticos para presentación en jornadas/congreso & ARS\, residual & Presentar resultados en evento nacional \\
\bottomrule
\end{longtable}

\textbf{Nota:} Dada la inflación en Argentina, se recomienda ejecutar las compras con prontitud para preservar el poder adquisitivo de los fondos asignados.

% ============================================================================
\section{Estado de entregables comprometidos}
% ============================================================================

\begin{longtable}{p{0.7cm}p{4.5cm}p{1.5cm}p{1.5cm}p{5cm}}
\toprule
\textbf{Cód.} & \textbf{Entregable} & \textbf{Estado} & \textbf{Avance} & \textbf{Observación} \\
\midrule
\multicolumn{5}{l}{\textit{Grupo A --- Técnico-científicos}} \\
\midrule
A1 & Marco normativo de calidad & \checkmark & 100\,\pct{} & Completado abr--may 2023 \\
A2 & Protocolo liofilización --- Ruibarbo & En proc. & $\sim$75\,\pct{} & Protocolo funcional; falta formalización final \\
A3 & Protocolo liofilización --- Salicornia & Canc. & --- & Matriz descartada \\
A4 & Protocolo liofilización --- Erizo & Canc. & --- & Matriz descartada \\
A5 & Resultados experimentales --- Ruibarbo & En proc. & $\sim$75\,\pct{} & 4 experimentos, 261 obs, análisis estadístico completo; faltan bioquímicos \\
A6 & Resultados experimentales --- Animal & Canc. & --- & Matriz descartada \\
\midrule
\multicolumn{5}{l}{\textit{Grupo B --- Transferencia al medio}} \\
\midrule
B1 & Manual técnico-económico & En proc. & $\sim$65\,\pct{} & Cap.~1--6 + Anexos A, C, D redactados; falta Anexo~B y revisión \\
B2 & Informe final del PID & Pend. & 0\,\pct{} & \textbf{Pendiente.} Este documento contribuye a B2 \\
B3 & Vinculación y transferencia & En proc. & $\sim$70\,\pct{} & Convenio Viamonte, presente en eventos \\
\midrule
\multicolumn{5}{l}{\textit{Grupo C --- Formación y difusión}} \\
\midrule
C1 & Consolidación grupo QA3DS & \checkmark & 100\,\pct{} & Grupo operativo, presencia digital, becas \\
C2 & Artículos de divulgación & En proc. & $\sim$80\,\pct{} & 2 artículos La Lupa redactados, pend. presentación \\
C3 & Presentaciones en jornadas & En proc. & $\sim$60\,\pct{} & Jornadas CyT 2023, Semana Ambiente Muni RG \\
\midrule
\multicolumn{5}{l}{\textit{Grupo D --- Emergentes}} \\
\midrule
D1 & Estudio de escalado & \checkmark & 100\,\pct{} & 5 documentos técnico-económicos \\
D2 & Artículo científico indexado & En proc. & $\sim$30\,\pct{} & Preprint completado; falta selec. revista y envío \\
D3 & Propuesta nuevo PID 2026 & Pend. & 0\,\pct{} & Objetivo del grupo: 1 PID/año \\
D4 & Dashboard integral & \checkmark & 100\,\pct{} & ETL + Dash, 4 pestañas, 12+ gráficos \\
D5 & Preprint científico LaTeX & \checkmark & 100\,\pct{} & 15 pp, 14 refs verificadas \\
\bottomrule
\end{longtable}

\textbf{Resumen:} De los 10 entregables activos (excluyendo 3 cancelados), 5 están completados (100\,\pct{}), 4 en proceso avanzado (60--80\,\pct{}) y 1 pendiente (B2, que es este informe).

% ============================================================================
\section{Alianzas estratégicas}
% ============================================================================

\begin{longtable}{p{4cm}p{2cm}p{7.5cm}}
\toprule
\textbf{Entidad} & \textbf{Estado} & \textbf{Descripción} \\
\midrule
Estancia Viamonte & Activo & Convenio de colaboración para materia prima (ruibarbo, ajo negro, grosellas) y entorno productivo. A 30\,km de Río Grande. \\
INTI & Referencia & Modelo de transferencia tecnológica (Trufas del Nuevo Mundo) \\
INTA & En gestión & Potencial colaboración técnica \\
IIByT-CONICET (UNC) & Puntual & Análisis DVS (sorption isotherms) \\
CIATI (Alto Valle) & Puntual & Análisis microbiológico y rotulado nutricional \\
ICTA (UNC) & Puntual & Asistencia técnica para evaluación sensorial \\
\bottomrule
\end{longtable}

% ============================================================================
\section{Dificultades y lecciones aprendidas}
% ============================================================================

\subsection{Dificultades encontradas}

\begin{enumerate}[nosep]
    \item \textbf{Reducción del alcance:} De las tres matrices originales (ruibarbo, salicornia, erizo de mar), solo se pudo trabajar experimentalmente con ruibarbo. Las otras fueron descartadas por limitaciones logísticas, de acceso a materia prima y de infraestructura.

    \item \textbf{Limitaciones del equipamiento:} El liofilizador RIFICOR LT-8 no siempre alcanzó las especificaciones del manual (temperatura y presión reales difirieron en algunos ciclos). La falta de mesa antivibratoria limita la precisión de las mediciones analíticas.

    \item \textbf{Rotación de alumnos:} 9 de los 16 alumnos egresaron o se retiraron durante el proyecto, lo que generó desafíos en la continuidad del trabajo experimental.

    \item \textbf{Distancia geográfica:} El equipo dividido entre Ushuaia y Río Grande añade complejidad logística. El laboratorio experimental está en Río Grande.

    \item \textbf{Contexto económico:} La inflación argentina dificulta la planificación presupuestaria y reduce el poder adquisitivo de los fondos asignados entre el momento de la aprobación y la ejecución.

    \item \textbf{Acceso a análisis externos:} Los laboratorios de análisis bioquímico y microbiológico se encuentran fuera de la provincia, lo que encarece y enlentece la obtención de resultados complementarios.
\end{enumerate}

\subsection{Lecciones aprendidas}

\begin{enumerate}[nosep]
    \item Concentrar esfuerzos en una matriz permitió obtener resultados más profundos y estadísticamente robustos que dispersar en tres líneas.
    \item La digitalización del flujo de trabajo (GitHub, notebook, dashboard) mejoró sustancialmente la trazabilidad y la productividad del equipo.
    \item El convenio con un productor real (Estancia Viamonte) resultó clave para el acceso sostenido a materia prima de calidad.
    \item La formación de alumnos en análisis de datos y estadística tiene impacto formativo que trasciende el proyecto.
    \item La generación de un estudio de escalado desde etapas tempranas orientó la discusión hacia la aplicabilidad real de los resultados.
\end{enumerate}

% ============================================================================
\section{Actividades pendientes y compromisos abiertos}
% ============================================================================

\begin{longtable}{p{0.5cm}p{5.5cm}p{2cm}p{5.5cm}}
\toprule
\textbf{\#} & \textbf{Actividad} & \textbf{Prioridad} & \textbf{Plazo estimado} \\
\midrule
1 & Formalizar B2 (informe final PIDPP) a partir de este documento & Crítica & Inmediato \\
2 & Completar B1 (Anexo B --- datos terceros, revisión docente) & Alta & 1--2 meses \\
3 & Enviar artículos La Lupa redactados para evaluación & Alta & 1 mes \\
4 & Ejecutar crédito presupuestario 2025 (ARS\,243.561) & Alta & Inmediato (preservar valor) \\
5 & Análisis bioquímicos externos (DVS, microbiología, rotulado) & Alta & 2--3 meses \\
6 & Seleccionar revista y enviar preprint como artículo (D2) & Media & 2--3 meses \\
7 & Completar análisis de actividad de agua (aw) & Media & 2--4 meses \\
8 & Formular propuesta de nuevo PID (D3) & Media & 3--6 meses \\
9 & Presentar en congreso/jornadas con resultados completos & Baja & 6 meses \\
10 & Avanzar con propuesta de planta piloto (Fase 1 del escalado) & Baja & 6--12 meses \\
\bottomrule
\end{longtable}

% ============================================================================
\section{Conclusiones}
% ============================================================================

El PID PP9884 ha producido resultados significativos en la caracterización del proceso de liofilización aplicado a ruibarbo patagónico, con un dataset robusto de 261 observaciones, análisis estadístico completo y un preprint científico listo para sometimiento a revisión por pares. Si bien el alcance original se redujo de tres matrices a una, la profundidad alcanzada en la matriz de ruibarbo supera lo inicialmente planificado en términos de análisis de datos.

El proyecto ha cumplido satisfactoriamente con los objetivos de consolidación del grupo QA3DS, formación de recursos humanos (16 alumnos, 6 becados), generación de know-how sobre liofilización y vinculación con el sector productivo regional.

Los principales compromisos abiertos son: la formalización del informe final (B2), el envío del artículo científico a revista indexada, y la ejecución del crédito presupuestario recientemente notificado.

Se recomienda que las actividades pendientes se integren en la propuesta del próximo PID del grupo QA3DS, asegurando la continuidad de las líneas de investigación iniciadas.

\vspace{1cm}

\noindent\rule{6cm}{0.5pt}\\
\textbf{Ing. Pesq. Ariel Luján Giamportone}\\
Director --- PID PP9884\\
Grupo QA3DS --- UTN FRTDF\\
Marzo 2026

\end{document}
